Tourist purchases do not determine where residents spend or where their income
changes. In the hypothetical city, 300 areas have large tourist-spending increases.
Of these, 110 have low resident-spending increases and 100 have high increases.
The contrast rules out a simple descriptive equivalence between the two local
spending measures. It does not establish that visitors displace residents or that
tourist demand causes either pattern.

The data distinguish these purchases through different links. A tourist's card
origin and purchase destination measure demand at that destination. Resident
records with an observed home address instead connect purchases to residential
origins, including purchases elsewhere. Workplace and ownership links provide
another route: a destination can pay workers who live in other areas, and those
workers can spend their income away from that destination. These links help
explain why local tourist purchases need not track local resident purchases.
Unobserved changes in amenities and common seasonal movements can also produce
the contrast. Recovering the effect of tourist demand therefore requires more
than adding the links to a map.

We estimate resident-spending responses using bilateral origin-destination
observations. The specification includes direct tourist exposure and exposure
through baseline work and purchase links, together with pair and month effects
and predictable direct and network controls. The baseline exposure is a constructed
regressor, not the equilibrium response that we seek to estimate. Variation comes
from changes in spending across tourist origin countries interacted with initial
destination composition. This variation helps distinguish changes in tourist
demand from resident purchases, but does not become exogenous merely because its
shares are measured before the shock. The identifying restriction excludes omitted
effects correlated with these constructed shocks after the controls. Correlated
seasonal changes or destination amenities would violate that restriction if the
controls did not remove their relevant variation.

Even under the identifying restriction, a resident-spending response is not an
income response. Its origin component combines changes in income with changes
in destination prices and attractiveness. A stated demand restriction permits
initial expenditure shares to correct for the destination component. When the
restriction is imposed through adjusted regressors, their estimated origin
coefficients describe the restricted income response; they are not unrestricted
origin-month fixed effects that separately measure income. This distinction also
limits the observed-income comparison. Month effects absorb common monthly
changes, so the comparison can examine spatial differences without independently
recovering the common aggregate income change. The remaining empirical task is
to assess the income-recovery restriction, not to treat a good spatial fit as
validation of the aggregate welfare mean.
